\SourceLeaf{052}
minimum unum reperiri debere ita comparatum. ut inter factores numeri, qui eius
ordinem exprimit, binarius saltem non pluries occurrat, quam inter factores numeri $m$, qui exprimit ordinem functionis $Y$: puta, si statuatur $m=k\, .\,2^\mu$, denotante $k$ numerum imparem, inter factores functionis $Y$ ad minimum unus reperietur ad ordinem $k'\, .\,2^\nu$ referendus, ita ut etiam $k'$ sit impar, atque vel $\nu=\mu$,
vel $\nu<\mu$. Veritas huius assertionis sponte sequitur inde, quod $m$ est aggregatum numerorum, qui ordinem singulorum factorum ipsius $Y$ exprimunt.

\SecondArticle{11}
Antequam ulterius progrediamur, expressionem quandam explicabimus, cuius introductio in omnibus de functionibus symmetricis disquisitionibus maximam
utilitatem affert, et quae nobis quoque peropportuna erit. Supponamus, $M$ esse
functionem quandam ex indeterminatis $a,b,c$ etc., et quidem sit $\mu$ multitudo
earum, quae in expressionem $M$ ingrediuntur, nullo respectu habito aliarum indeterminatarum, si quas forte implicet ipsa $M$. Permutatis illis $\mu$ indeterminatis omnibus quibus fieri potest modis tum inter se tum cum $m-\mu$ reliquis ex
$a,b,c$ etc., orientur ex $M$ aliae expressiones ipsi $M$ similes, ita ut omnino habeantur
\[
m(m-1)(m-2)(m-3)\, .\, .\, .\, .\, .\,(m-\mu+1)
\]
expressiones, ipsa $M$ inclusa, quarum complexum simpliciter dicemus \textit{complexum
omnium $M$.} Hinc sponte patet, quid significet aggregatum omnium $M$, productum ex omnibus $M$ etc. Ita e. g. $\pi$ dicetur productum ex omnibus $a-b$,
$\upsilon$ productum ex omnibus $x-a$, $\upsilon'$ aggregatum omnium $\frac{\upsilon}{x-a}$ etc.

Si forte $M$ est functio symmetrica respectu quarundam ex $\mu$ indeterminatis, quas continet, istarum permutationes inter se functionem $M$ non variant, quamobrem in complexu omnium $M$ quilibet terminus pluries, et quidem
$1\, .\,2\, .\,3\, .\, .\, .\, .\, .\,\nu$ vicibus reperietur, si $\nu$ est multitudo indeterminatarum, quarum
respectu $M$ est symmetrica. Si vero $M$ non solum respectu $\nu$ indeterminatarum
symmetrica est, sed insuper respectu $\nu'$ aliarum, nec non respectu $\nu''$ aliarum etc.,
ipsa $M$ non variabitur sive binae e primis $\nu$ indeterminatis inter se permutentur,
sive binae e secundis $\nu'$, sive binae e tertiis $\nu''$ etc., ita ut semper
\[
1\, .\,2\, .\,3\, .\, .\, .\, .\, .\,\nu\, .\quad1\, .\,2\, .\,3\, .\, .\, .\, .\, .\,\nu'\, .\quad1\, .\,2\, .\,3\, .\, .\, .\, .\, .\,\nu''\text{ etc.}
\]
\SourceLeaf{053}
permutationes terminis identicis respondeant. Quare si ex his terminis identicis
semper unicum tantum retineamus, omnino habebimus
\[
\frac{m(m-1)(m-2)(m-3)\, .\, .\, .\, .\, .\,(m-\mu+1)}{1\, .\,2\, .\,3\, .\, .\, .\, .\, .\,\nu\, .\,1\, .\,2\, .\,3\, .\, .\, .\, .\, .\,\nu'\, .\,1\, .\,2\, .\,3\, .\, .\, .\, .\, .\,\nu''\text{ etc.}}
\]
terminos, quorum complexum dicemus \textit{complexum omnium $M$ exclusis repetitionibus}, ut a \textit{complexu omnium $M$ admissis repetitionibus} distinguatur. Quoties nihil
expressis verbis monitum fuerit, repetitiones admitti semper subintelligemus.

Ceterum facile perspicietur, aggregatum omnium $M$, vel productum ex
omnibus $M$, vel generaliter quamlibet functionem symmetricam omnium $M$ semper fieri functionem .symmetricam indeterminatarum $a,b,c$ etc., sive admittantur repetitiones, sive excludantur.

\SecondArticle{12}
Iam considerabimus, denotantibus $u,x$ indeterminatas, productum ex omnibus $u-(a+b)x+ab$, exclusis repetitionibus, quod per $\zeta$ designabimus. Erit
itaque $\zeta$ productum ex $\frac12m(m-1)$ factoribus his
\[
\begin{gathered}
u-(a+b)x+ab\\
u-(a+c)x+ac\\
u-(a+d)x+ad\\
\text{etc.}\\
u-(b+c)x+bc\\
u-(b+d)x+bd\\
\text{etc.}\\
u-(c+d)x+cd\\
\text{etc. etc.}
\end{gathered}
\]
Quae functio quum indeterminatas $a,b,c$ etc. symmetrice implicet, assignari
poterit functio integra indeterminatarum $u,x,l',l'',l'''$ etc., per $z$ denotanda,
quae transeat in $\zeta$, si loco indeterminatarum $l',l'',l'''$ etc. substituantur $\lambda',\lambda'',\lambda'''$ etc.
Denique designemus per $Z$ functionem solarum indeterminatarum $u,x$, in
quam $z$ transit, si indeterminatis $l',l'',l'''$ etc. tribuamus valores determinatos
$L',L'',L'''$ etc.

Hae tres functiones $\zeta,z,Z$ considerari possunt tamquam functiones integrae ordinis $\frac12m(m-1)$ indeterminatae $u$ cum coëfficientibus indeterminatis, qui
\SourceLeaf{054}
quidem coëfficientes erunt
\[
\begin{aligned}
&\text{pro }\zeta,\quad\text{functiones indeterminatarum }x,a,b,c\text{ etc.}\\
&\text{pro }z,\quad\text{functiones indeterminatarum }x,l',l'',l'''\text{ etc.}\\
&\text{pro }Z,\quad\text{functiones solius indeterminatae }x.
\end{aligned}
\]
Singuli vero coëfficientes ipsius $z$ transibunt in coëfficientes ipsius $\zeta$ per substitutiones $l'=\lambda',l''=\lambda'',l'''=\lambda'''$ etc. nec non in coëfficientes ipsius $Z$ per substitutiones $l'=L',l''=L'',l'''=L'''$ etc. Eadem, quae modo de coëfficientibus diximus, etiam de determinantibus functionum $\zeta,z,Z$ valebunt. Atque in
hos ipsos iam propius inquiremus, et quidem eum in finem, ut demonstretur

\textsc{Theorema.} \textit{Quoties non est $P=0$, determinans functionis $Z$ certo nequit
esse identice $=0$.}

\SecondArticle{13}
Perfacilis quidem esset demonstratio huius theorematis. si supponere liceret, $Y$ resolvi posse in factores simplices
\[
(x-A)(x-B)(x-C)(x-D)\, .\, .\, .\, .\, .
\]
Tunc enim certum quoque esset, $Z$ esse productum ex omnibus $u-(A+B)x+AB$,
atque determinantem functionis $Z$ productum e differentiis inter binas quantitatum
\[
\begin{gathered}
(A+B)x-AB\\
(A+C)x-AC\\
(A+D)x-AD\\
\text{etc.}\\
(B+C)x-BC\\
(B+D)x-BD\\
\text{etc.}\\
(C+D)x-CD\\
\text{etc. etc}
\end{gathered}
\]
Hoc vero productum identice evanescere nequit, nisi aliquis factorum per se identice fiat $=0$, unde sequeretur, duas quantitatum $A,B,C$ etc. aequales esse,
adeoque determinantem $P$ functionis $Y$ fieri $=0$, contra hyp.

\SourceLeaf{055}
At seposita tali argumentatione, quam ad instar art. 6 a petitione principii
proficisci manifestum est, statim ad demonstrationem stabilem theorematis art. 12
explicandam progredimur.

\SecondArticle{14}
Determinans functionis $\zeta$ erit productum ex omnibus differentiis inter binas $(a+b)x-ab$, quarum differentiarum multitudo est
\[
\frac12m(m-1)(\tfrac12m(m-1)-1)=\frac14(m+1)m(m-1)(m-2)
\]
Hic numerus itaque indicat ordinem determinantis functionis $\zeta$ respectu indeterminatae $x$. Determinans functionis $z$ quidem ad eundem ordinem pertinebit:
contra determinans functionis $Z$ utique ad ordinem inferiorem pertinere potest,
quoties scilicet quidam coëfficientes inde ab altissima potestate ipsius $x$ evanescunt.
Nostrum iam est demonstrare, in determinante functionis $Z$ \textit{omnes} certo coëfficientes evanescere non posse.

Propius considerando differentias illas, quarum productum est determinans functionis $\zeta$, deprehendemus, partem ex ipsis (puta differentias inter binas
$(a+b)x-ab$ tales, quae elementum commune habent) suppeditare
\[
\text{productum ex omnibus}\quad(a-b)(x-c)
\]
e reliquis vero (puta e differentiis inter binas $(a+b)x-ab$ tales, quarum elementa diversa sunt) oriri
\[
\text{productum ex omnibus}\quad(a+b-c-d)x-ab+cd,\quad\text{exclusis repetitionibus.}
\]
Productum prius factorem unumquemque $a-b$ manifesto $m-2$ vicibus continebit, quemvis factorem $x-c$ autem $(m-1)(m-2)$ vicibus, unde facile concludimus, hocce productum fieri
\[
=\pi^{m-2}\upsilon^{(m-1)(m-2)}
\]
Quodsi ita productum posterius per $\rho$ designamus, determinans functionis $\zeta$ erit
\[
=\pi^{m-2}\upsilon^{(m-1)(m-2)}\rho
\]
Denotando porro per $r$ functionem indeterminatarum $x,l',l'',l'''$ etc. eam, quae
transit in $\rho$ per substitutiones $l'=\lambda',l''=\lambda'',l'''=\lambda'''$ etc., nec non per $R$
